Deterministic backoff with listen before talk (DB-LBT) uses local countdown
interruptions to build collision-reduced Wi-Fi/NR-U access schedules. Its
recovery parameters are normally fixed, although offered load, active
contenders, and external occupancy change the backoff profile that works best.
We formulate DB-LBT parameter adaptation as a protocol-constrained contextual
bandit problem under locally observable MAC information. Each transmitter runs
LinUCB over 24 legal $(\kappa,\beta,m,b_{\mathrm{init}})$ tuples. The context
contains only local CCA, queue, delay, retransmission, and ACK/HARQ statistics;
the selected action changes the next DB-LBT recovery counter but leaves native
Wi-Fi and NR-U procedures intact.

We derive the sensing, reception, separability, and timescale conditions that
make this local context informative. Paired event experiments then compare
adaptive DB-LBT with random LBT and fixed DB-LBT. In a static $4+4$ regime,
adaptive and fixed TMC DB-LBT are numerically identical. Across held-out
scenarios, adaptation raises mean utility from $0.926467$ to $0.928023$; the
non-ideal gain is $0.289\%$. LinUCB also exceeds context-free UCB in the
combined dynamic case. A packet-level ns-3 cross-check produces mixed
technology-level directions and exposes the unresolved validation gap.
